\par\noindent\textbf{(1)} 面密度涨落 \(\delta n\) 产生面电荷 \(\delta\sigma=-e\delta n\)。对单一横向波数 \(k\)，平面两侧解Laplace方程得 \(\phi(z)=\phi_0e^{-k|z|}\)。法向电位移跃变 \(\epsilon(E_z^+-E_z^-)=\delta\sigma\)，而 \(E_z^+=k\phi_0\)、\(E_z^-=-k\phi_0\)，因此 \(\phi_0=-e\delta n/(2\epsilon k)\)。
\par\noindent\textbf{(2)} 二维电子连续方程 \(-i\omega\delta n+ikn_sv=0\)，动量式 \(-i\omega mv=-eE_x\)，其中 \(E_x=-ik\phi_0\)。代入上一问的 \(\phi_0\) 消元，得到 \(\omega^2=n_se^2k/(2m\epsilon)\)。二维等离激元是 \(\omega\propto\sqrt k\) 的非局域静电模式，不等同于三维体等离子体具有非零的 \(k=0\) 等离子频隙。
\par\noindent\textbf{(3)} 记 \(C=\sqrt{n_se^2/(2m\epsilon)}\)，则 \(\omega=C\sqrt k\)，相速度 \(v_p=\omega/k=C/\sqrt k\)，群速度 \(v_g=d\omega/dk=C/(2\sqrt k)=v_p/2\)。形式上 \(k\to0\) 的相群速度都发散，说明准静态无迟滞模型不能无限外推；在电磁迟滞尺度需改用二维等离激元—极化激元。
\par\noindent\textbf{(4)} 若栅极在距离 \(d\) 处接地，像电荷改变二维Fourier库仑核为 \(\phi_0=(\delta\sigma/(2\epsilon k))(1-e^{-2kd})\)。因此 \(\omega^2=[n_se^2k/(2m\epsilon)](1-e^{-2kd})\)。当 \(kd\ll1\) 时 \(1-e^{-2kd}\simeq2kd\)，得到声学模 \(\omega\simeq ck\)，\(c=\sqrt{n_se^2d/(m\epsilon)}\)，金属栅屏蔽改变了长波指数。
